\section{Haces cuasicoherentes y haces coherentes}
\label{section:0.5-fr}

\subsection{Haces cuasicoherentes}
\label{subsection:0.5.1-fr}

\begin{env}[5.1.1]
\label{0.5.1.1-fr}
Sean $(X,\mathcal{O}_X)$ un espacio anillado, $\mathcal{F}$ un
$\mathcal{O}_X$-Módulo. El dato de un homomorfismo
$u:\mathcal{O}_X\to\mathcal{F}$ de $\mathcal{O}_X$-Módulos equivale al
de la sección $s=u(1)\in\Gamma(X,\mathcal{F})$. En efecto, cuando $s$ está
dada, para toda sección $t\in\Gamma(U,\mathcal{O}_X)$, se tiene necesariamente
$u(t)=t\cdot(s|U)$; se dice que $u$ está \emph{definido por la sección $s$}. Si
ahora $I$ es un conjunto cualquiera de índices, consideremos el haz
suma directa $\mathcal{O}_X^{(I)}$, y para todo $i\in I$, sea $h_i$
la inyección canónica del $i$-ésimo factor en $\mathcal{O}_X^{(I)}$; se sabe
que $u\to(u\circ h_i)$ es un isomorfismo de
$\operatorname{Hom}_{\mathcal{O}_X}(\mathcal{O}_X^{(I)},\mathcal{F})$ sobre el
producto
$\bigl(\operatorname{Hom}_{\mathcal{O}_X}(\mathcal{O}_X,\mathcal{F})\bigr)^I$.
Hay, pues, una correspondencia biunívoca canónica entre los homomorfismos
$u:\mathcal{O}_X^{(I)}\to\mathcal{F}$ y las \emph{familias de secciones
$(s_i)_{i\in I}$ de $\mathcal{F}$ sobre $X$}. El homomorfismo $u$
correspondiente a $(s_i)$ aplica un elemento
$(a_i)\in(\Gamma(U,\mathcal{O}_X))^{(I)}$ en
$\sum_{i\in I}a_i\cdot(s_i|U)$.

Se dice que $\mathcal{F}$ está \emph{generado por la familia $(s_i)$} si
el homomorfismo $\mathcal{O}_X^{(I)}\to\mathcal{F}$ definido
\oldpage[0\textsubscript{I}]{45}
por esta familia es \emph{sobreyectivo} (es decir, si, para todo
$x\in X$, $\mathcal{F}_x$ es un $\mathcal{O}_x$-módulo generado por los
$(s_i)_x$). Se dice que $\mathcal{F}$ está \emph{generado por sus secciones
sobre $X$} si está generado por la familia de todas estas secciones (o
por una subfamilia), es decir, si existe un homomorfismo sobreyectivo
$\mathcal{O}_X^{(I)}\to\mathcal{F}$ para un $I$ adecuado.

Se observará que un $\mathcal{O}_X$-Módulo $\mathcal{F}$ puede ser tal que
exista un punto $x_0\in X$ para el cual $\mathcal{F}|U$ no esté generado por
sus secciones sobre $U$, \emph{cualquiera que sea la vecindad abierta $U$ de
$x_0$}: basta tomar $X=\mathbb{R}$, como $\mathcal{O}_X$ el haz
simple $\mathbb{Z}$, como $\mathcal{F}$ el subhaz algebraico de
$\mathcal{O}_X$ tal que $\mathcal{F}_0=\{0\}$, $\mathcal{F}_x=\mathbb{Z}$ para
$x\ne0$, y, finalmente, $x_0=0$: la única sección de $\mathcal{F}|U$ sobre
$U$ es $0$ para una vecindad $U$ de $0$.
\end{env}

\begin{env}[5.1.2]
\label{0.5.1.2-fr}
Sea $f:X\to Y$ un morfismo de espacios anillados. Si $\mathcal{F}$ es un
$\mathcal{O}_X$-Módulo generado por sus secciones sobre $X$, entonces
el homomorfismo canónico
$f^*(f_*(\mathcal{F}))\to\mathcal{F}$ (4.4.3.3) es \emph{sobreyectivo}: en
efecto, con las notaciones de (5.1.1), $s_i\otimes1$ es una sección de
$f^*(f_*(\mathcal{F}))$ sobre $X$, y su imagen en $\mathcal{F}$ es
$s_i$. El ejemplo de (5.1.1), en el que $f$ es la identidad, muestra que el recíproco de
esta proposición es inexacto en general.
\end{env}

\begin{env}[5.1.3]
\label{0.5.1.3-fr}
Se dice que un $\mathcal{O}_X$-Módulo $\mathcal{F}$ es \emph{cuasicoherente}
si, para todo $x\in X$, hay una vecindad abierta $U$ de $x$ tal que
$\mathcal{F}|U$ sea isomorfo al \emph{conúcleo} de un homomorfismo de la
forma
$\mathcal{O}_X^{(I)}|U\to\mathcal{O}_X^{(J)}|U$, donde $I$ y $J$ son
conjuntos arbitrarios de índices. Es claro que el propio $\mathcal{O}_X$ es
un $\mathcal{O}_X$-Módulo cuasicoherente, y que toda suma directa finita de
$\mathcal{O}_X$-Módulos cuasicoherentes es un $\mathcal{O}_X$-Módulo
cuasicoherente. Se dice que una $\mathcal{O}_X$-\emph{Álgebra} $\mathcal{A}$ es
\emph{cuasicoherente} si lo es como $\mathcal{O}_X$-Módulo.
\end{env}

\begin{env}[5.1.4]
\label{0.5.1.4-fr}
Sea $f:X\to Y$ un morfismo de espacios anillados. Si $\mathcal{G}$ es un
$\mathcal{O}_Y$-Módulo cuasicoherente, entonces $f^*(\mathcal{G})$ es un
$\mathcal{O}_X$-Módulo cuasicoherente. En efecto, para todo $x\in X$, hay
una vecindad abierta $V$ de $f(x)$ en $Y$ tal que $\mathcal{G}|V$ es el
conúcleo de un homomorfismo
$\mathcal{O}_Y^{(I)}|V\to\mathcal{O}_Y^{(J)}|V$. Si $U=f^{-1}(V)$, y si
$f_U$ es la restricción de $f$ a $U$, se tiene
$f^*(\mathcal{G})|U=f_U^*(\mathcal{G}|V)$; como $f_U^*$ es exacto por la derecha
y conmuta con las sumas directas, $f_U^*(\mathcal{G}|V)$ es conúcleo de un
homomorfismo
$\mathcal{O}_X^{(I)}|U\to\mathcal{O}_X^{(J)}|U$.
\end{env}

